
\clearpage
\part*{Part VII --- The other Shor: elliptic-curve discrete logarithms}
\addcontentsline{toc}{part}{Part VII --- The other Shor: elliptic-curve discrete logarithms}

\begin{keybox}[title={Why a second half}]
Everything so far factors integers. The same three layers --- Hadamards, an oracle ladder, an
inverse QFT --- also break elliptic-curve cryptography, and the outer two layers are unchanged. What
changes is the middle, and it changes completely.

For factoring, the group is $(\mathbb{Z}/N)^\times$ and the oracle multiplies by a
\emph{classical} constant. Multiplication by a constant is cheap: Part~II built it out of shifted
additions. For the discrete logarithm on a curve, the group is $E(\mathbb{F}_p)$ and the oracle adds
a classical \emph{point} --- and the chord-and-tangent law needs a division, i.e.\ a modular
\emph{inverse}, of a quantum value. That single fact is the whole story of this part: an inversion
costs an order of magnitude more than the multiplication it replaces, and every optimisation below
is an attempt to make it cheaper, do it less often, or avoid it altogether.
\end{keybox}

\section{What changes, and what does not}\label{ec:sec:orient}

Let $E$ be an elliptic curve over $\mathbb{F}_p$, $P$ a point of prime order $r$, and
$Q=[k]P$ for an unknown $k$. Recovering $k$ is the \emph{elliptic-curve discrete logarithm problem}
(ECDLP), and it is what ECDSA, ECDH and every deployed curve-based scheme rest on.

Shor's algorithm for it \cite{ec:proos-zalka} uses two counting registers instead of one:
\[
\frac{1}{2^{m}}\sum_{u,v}\ket u\ket v\ket{\mathcal O}
\;\xrightarrow{\ \text{oracle}\ }\;
\frac{1}{2^{m}}\sum_{u,v}\ket u\ket v\ket{[u]P+[v]Q}
\;\xrightarrow{\ F^\dagger\otimes F^\dagger\ }\;
\text{measure } (j_1,j_2).
\]
The third register depends on $(u,v)$ only through $u+kv \bmod r$, so its level sets are cosets of
the lattice generated by $(-k,1)$ and $(r,0)$. The Fourier transform concentrates on the dual
lattice, and a measured pair obeys, with $a_i=\mathrm{round}(j_i r/2^{m})$,
\begin{equation}\label{ec:eq:post}
a_2 \equiv k\,a_1 \pmod r,
\qquad\text{so}\qquad
k = a_2\,a_1^{-1} \bmod r
\end{equation}
whenever $a_1$ is invertible --- which discards $j_1=0$ automatically. One measurement suffices with
constant probability; \S\ref{ec:sec:verify} shows the distribution this produces on a toy curve, and
\S\ref{ec:sec:post} the search that \cite{ec:ekera} adds when $2^m$ is not a multiple of $r$.

\begin{center}
\begin{quantikz}[column sep=8pt,row sep=8pt]
\lstick{$\ket{u_0}$}     & \gate{H} & \ctrl{6} & \qw          & \ \ldots\ & \qw                & \qw       & \ \ldots\ & \qw                & \gate[3]{F^\dagger} & \meter{}\\
\lstick{$\vdots$}        &          &          &              &           &                    &           &           &                    &                     & \\
\lstick{$\ket{u_{m-1}}$} & \gate{H} & \qw      & \qw          & \ \ldots\ & \ctrl{4}           & \qw       & \ \ldots\ & \qw                &                     & \meter{}\\
\lstick{$\ket{v_0}$}     & \gate{H} & \qw      & \qw          & \ \ldots\ & \qw                & \ctrl{3}  & \ \ldots\ & \qw                & \gate[3]{F^\dagger} & \meter{}\\
\lstick{$\vdots$}        &          &          &              &           &                    &           &           &                    &                     & \\
\lstick{$\ket{v_{m-1}}$} & \gate{H} & \qw      & \qw          & \ \ldots\ & \qw                & \qw       & \ \ldots\ & \ctrl{1}           &                     & \meter{}\\
\lstick{$\ket{S}$}       & \qw      & \gate{+P}& \gate{+[2]P} & \ \ldots\ & \gate{+[2^{m-1}]P} & \gate{+Q} & \ \ldots\ & \gate{+[2^{m-1}]Q} & \qw                 & \qw
\end{quantikz}
\end{center}

\noindent
Each rung adds a \emph{classical} point --- the multiples $[2^i]P$ and $[2^i]Q$ are precomputed on a
laptop --- controlled by one counting qubit. There are $2m\approx 2n$ of them, against the $2n$
rungs of Part~II's ladder, so at the level of \emph{counting} nothing has changed. The cost is
entirely inside the box marked $+P$.

\begin{warnbox}[title={The accumulator never starts at infinity}]
The diagram starts the third register at a fixed point $S$, not at the identity $\mathcal O$. Affine
addition formulas cannot represent the point at infinity at all: there is no $(x,y)$ for it, and the
slope $\lambda$ is undefined the moment the accumulator equals $\pm$ the addend. Starting at a
constant $S$ and reading the answer relative to it keeps the accumulator on an ordinary point
throughout, and shifts the computed coset without touching the interference. The same device
appears as the offset $T$ in the windowed constructions of \S\ref{ec:sec:window}, and as
\code{P0} in the Classiq tutorial \cite{ec:classiq}.
\end{warnbox}

\section{Why it works}\label{ec:sec:why}

Section~\ref{ec:sec:orient} asserted \eqref{ec:eq:post}. This section derives it. The argument is
short, it is the only place in this part where the \emph{quantum} content lives --- everything after
\S\ref{ec:sec:grouplaw} is reversible arithmetic --- and it is checked numerically in
\code{bench/ec\_ablation.py}, which evaluates the amplitudes by brute force and compares them
against the closed form below.

\subsection{The state, in four lines}

Write $q=2^m$ for the size of each counting register and $\omega=e^{2\pi i/q}$. Start from a uniform
superposition and apply the oracle:
\[
\frac1q\sum_{u,v}\ket u\ket v\ket S
\;\longrightarrow\;
\frac1q\sum_{u,v}\ket u\ket v\,\big|\,S+[u]P+[v]Q\,\big\rangle
=\frac1q\sum_{u,v}\ket u\ket v\,\big|\,S+[u+kv]P\,\big\rangle ,
\]
the last equality being the only place $Q=[k]P$ is used. The third register therefore depends on
$(u,v)$ \emph{only} through
\[
w \;=\; u+kv \bmod r ,
\]
so group the terms by $w$:
\begin{equation}\label{ec:eq:grouped}
\frac1q\sum_{w\in\mathbb{Z}_r}\Bigg(\sum_{(u,v)\,:\,u+kv\equiv w}\ket u\ket v\Bigg)\big|\,S+[w]P\,\big\rangle .
\end{equation}
The third register is never measured and never needs to be: \eqref{ec:eq:grouped} is already a sum
of orthogonal branches, so the reduced state on $(u,v)$ is a mixture over $w$, and it is enough to
analyse one branch.

\begin{keybox}[title={Why the offset $S$ cannot matter}]
$S$ appears only inside the third register's label. It relabels which basis state each branch
attaches to; it does not touch the $(u,v)$ amplitudes at all, so it cannot change the measured
distribution. The same argument covers $w$: every branch of \eqref{ec:eq:grouped} gives the same
$|$amplitude$|$, which the benchmark confirms directly by computing two different branches and
comparing (result: \code{@@deriv_offset@@}). $S$ is there for the arithmetic --- to keep the
accumulator off the point at infinity --- and for nothing else.
\end{keybox}

\subsection{Two lattices}

Fix a branch $w$. Its $(u,v)$ set is a coset of
\[
L \;=\; \{(u,v)\in\mathbb{Z}^2 \;:\; u+kv\equiv 0 \ (\mathrm{mod}\ r)\}
\;=\; \big\langle\, (r,0),\ (-k,1) \,\big\rangle ,
\]
the \emph{period lattice}. The two generators say the two things that can be done without changing
the answer: take $r$ steps in $u$, or take one step in $v$ and $k$ steps back in $u$. Its
determinant is $r$, so $[\mathbb{Z}^2 : L]=r$ and the plane splits into exactly $r$ cosets --- one
per value of $w$, as it must.

\emph{The discrete logarithm is the slope of the second generator.} Finding $L$ is finding $k$;
that is the whole problem, restated.

A Fourier transform does not see $L$, it sees the \emph{dual}
\[
L^{*} \;=\; \{\,y\in\mathbb{Q}^2 \;:\; \langle y,x\rangle\in\mathbb{Z}\ \ \forall x\in L\,\}
\;=\; \big\langle\, \tfrac1r(1,k),\ (0,1) \,\big\rangle ,
\]
obtained as $(B^{-1})^{\mathsf T}$ from $L$'s basis $B$. Writing a point of $L^{*}$ as
$(a_1/r,\,a_2/r)$ and clearing denominators, membership says exactly
\[
a_2 \equiv k\,a_1 \ (\mathrm{mod}\ r) ,
\]
which is \eqref{ec:eq:post}. It is not a consequence of the dual lattice --- it \emph{is} the dual
lattice, written in coordinates.

\subsection{The amplitude}

Take $r\mid q$ for a moment and write $q=rM$. Then for each $v$ the admissible $u$ in $[0,q)$ are
$u_0(v)+tr$ with $u_0(v)=(w-kv)\bmod r$ and $t=0,\dots,M-1$ --- exactly $M$ of them. The amplitude
at $(j_1,j_2)$ after the transform is, up to normalisation,
\[
A(j_1,j_2)\;\propto\;\sum_{v=0}^{q-1}\omega^{v j_2}\,\omega^{u_0(v)j_1}
\underbrace{\sum_{t=0}^{M-1}\omega^{t r j_1}}_{\text{(i)}} .
\]
Sum (i) is geometric with ratio $\omega^{rj_1}$: it equals $M$ when $rj_1\equiv 0 \pmod q$, i.e.\
when $M \mid j_1$, and \emph{exactly zero} otherwise. So only $j_1=a_1M$ survives. For those,
$\omega^{u_0 j_1}=e^{2\pi i\,u_0 a_1/r}$ depends on $u_0$ only mod $r$ --- which is why reducing
$u_0$ into $[0,r)$ above was harmless --- and $u_0\equiv w-kv$, leaving
\[
A \;\propto\; e^{2\pi i\,w a_1/r}\sum_{v=0}^{q-1} e^{2\pi i\,v\,(j_2-k a_1 M)/q} .
\]
That second sum is $q$ when $j_2\equiv ka_1M \pmod q$ and zero otherwise. Writing $j_2=a_2M$ and
dividing through by $M$:
\[
\boxed{\;a_2\equiv k\,a_1 \pmod r\;}
\]
and every surviving outcome has $|A|^2=1/r$. The distribution is \emph{uniform on the $r$ points of
the dual lattice}, and nothing else has any amplitude at all.

\begin{center}\small
\begin{tabular}{@{}rrrrrcr@{}}
\toprule
$r$ & $k$ & $q$ & outcomes & each & closed form exact & usable $\varphi(r)/r$\\
\midrule
%%DERIVCLEAN%%
\bottomrule
\end{tabular}
\end{center}

\subsection{Reading off $k$, and what it costs}

Given $(a_1,a_2)$ with $a_2\equiv ka_1$, recover $k=a_2a_1^{-1}\bmod r$ --- possible exactly when
$a_1$ is invertible. For prime $r$ that means $a_1\neq0$, so $r-1$ of the $r$ outcomes work and one
run succeeds with probability $(r-1)/r$: overwhelming, and it does not improve with more qubits
because the $a_1=0$ outcome is a fixed $1/r$ of the mass. Note $r$ is not secret --- it is the
curve's published group order --- so this is a division, not a search.

\begin{warnbox}[title={The clean case never actually happens}]
Everything above assumed $r\mid q$. For ECDLP $r$ is \emph{prime}, and $q$ is a power of two, so
$r\mid q$ only when $r=2$. The exact calculation is a device for seeing the structure; the real
situation is always the other one.

What goes wrong is only sum (i): with $r\nmid q$ it no longer vanishes, it becomes a Dirichlet
kernel. Amplitude leaks onto the integers either side of $a_1q/r$, which is why the measured peaks
of \S\ref{ec:sec:post} sit at $(6,5)$ and $(2,3)$ rather than at the exact dual-lattice points
$(1.6,3.2)$ and $(3.2,6.4)$, and why $a_i=\mathrm{round}(j_ir/q)$ has a rounding in it at all. The
leak shrinks as $q$ grows past $r$. Entries are $P(\text{usable and correct})$, i.e.\ the chance
of an outcome with $a_1$ invertible \emph{and} $a_2\equiv ka_1$; the last column is the irreducible
$a_1=0$:

\begin{center}\small
\begin{tabular}{@{}rrrrrr@{}}
\toprule
& \multicolumn{4}{c}{$m-\lceil\log_2 r\rceil$} & \\
\cmidrule(lr){2-5}
$r$ & $0$ & $1$ & $2$ & $3$ & $P(a_1{=}0)$\\
\midrule
%%DERIVREAL%%
\bottomrule
\end{tabular}
\end{center}

\noindent
At $m=\lceil\log_2 r\rceil$ the answer is barely better than a coin; three extra qubits per register
take it to $@@p13_m7@@\%$, with the residual $@@z13@@$ being the irreducible $a_1=0$ outcome. This
is the whole content of \cite{ec:ekera}'s analysis for the discrete logarithm, and the reason the
post-processing of \S\ref{ec:sec:post} searches a small neighbourhood of the rounded value rather
than trusting it.
\end{warnbox}

\section{The group law, and why it resists reversibility}\label{ec:sec:grouplaw}

For $P_1=(x_1,y_1)$ and $P_2=(x_2,y_2)$ on $y^2=x^3+ax+b$ with $x_1\neq x_2$,
\begin{equation}\label{ec:eq:add}
\lambda=\frac{y_1-y_2}{x_1-x_2},\qquad
x_3=\lambda^2-x_1-x_2,\qquad
y_3=\lambda(x_2-x_3)-y_2 .
\end{equation}
Three subtractions, one squaring, one multiplication --- and one division. In Part~II every
arithmetic step was addition or multiplication by a compile-time constant. Here the divisor
$x_1-x_2$ is a quantum value, so the inverse must be computed \emph{by circuit}, on a
superposition, and then uncomputed. Section~\ref{ec:sec:kaliski} does that, and
\S\ref{ec:sec:cost} shows it dominating everything else by a factor of three to five.

\begin{center}\small\setlength{\tabcolsep}{5pt}
\begin{tabular}{@{}rllr@{}}
\toprule
Level & Object & Made of & In code\\
\midrule
5 & the two ladders & $2m$ rungs, one per counting qubit & \code{ec\_shor.py}\\
4 & $c$-$(+P_i)$ & 4 sub, 1 add, 2 div, 2 mul, 1 sqr, 1 neg & \code{ec\_pointadd.py}\\
3 & $x^{-1}\bmod p$ & $2n$ Kaliski rounds + correction & \code{ec\_kaliski.py}\\
3$'$ & $x\cdot y\bmod p$ & $n$ modular doublings + $n$ additions & \code{ec\_mult.py}\\
2 & $\pm\bmod p$, $2x\bmod p$ & plain add/sub plus one reduction & \code{ec\_modarith.py}\\
1 & $\ADD$ & CDKM chain, or Gidney's AND chain & \code{ec\_adders.py}\\
0 & AND & $4T$ to compute, $0T$ to uncompute & \code{ec\_gates.py}\\
\bottomrule
\end{tabular}
\end{center}

\noindent
Level~3 is new and it is the expensive one. Everything below it is Part~II's tower with the
modulus renamed.

\section{Level 2: modular arithmetic over $\mathbb{F}_p$}\label{ec:sec:modarith}

Two blocks carry everything above them, and both are from Roetteler, Naehrig, Svore and Lauter
\cite{ec:rnsl17}. Neither is hard; both are worth seeing, because \S\ref{ec:sec:approx} makes its
savings by attacking exactly the parts described here as ``the awkward bit''.

\paragraph{Modular doubling.}
Shift, subtract $p$, add it back if that went negative:
\[
2x \;\longrightarrow\; 2x-p \;\longrightarrow\; \begin{cases} 2x-p & 2x\ge p\\ 2x & \text{otherwise.}\end{cases}
\]
The shift is free --- it renames wires. The awkward bit is the borrow flag, which must return to
$\ket0$ or it is garbage in the sense of \S\ref{sec:garbage}. It clears for one CNOT: $2x$ is even
and $p$ is odd, so \emph{the answer is even exactly when no reduction happened}, and the answer's
own low bit is the flag.

\paragraph{Modular addition.} The same shape, but the flag-clearing trick is prettier.
After adding and subtracting $p$, the flag holds $[\,x+y<p\,]$, i.e.\ ``no reduction''. And
\begin{itemize}\setlength\itemsep{1pt}
\item no reduction means the answer is $x+y$, which is $\ge x$;
\item a reduction means the answer is $x+y-p$, which is $<x$ because $y<p$.
\end{itemize}
So $[\,\text{answer}<x\,]$ is exactly the complement of the flag, and one comparison against an
operand that is still lying there clears it. No extra ancilla, no second copy. In code:

%%LST ec_modarith.cmodadd doc=1%%

\noindent
The doubling is the same shape and its flag-clearing is one CNOT --- the last two lines below,
where \code{xe[0]} is the answer's own low bit:

%%LST ec_modarith.moddbl doc=1%%

\section{Level 3: inversion by Kaliski's algorithm}\label{ec:sec:kaliski}

The binary extended Euclidean algorithm branches four ways per step: $u$ even, $v$ even, both odd
with $u>v$, both odd with $u\le v$. A quantum circuit cannot branch --- every branch must run on
every input, on a superposition. Kaliski's formulation as reworked by H\"aner, Jaques, Naehrig,
Roetteler and Soeken \cite{ec:hjnrs20} makes all four cases the \emph{same} straight line:

\begin{center}\small
\begin{minipage}{0.86\textwidth}
\begin{tabbing}
xxx\=xxxx\=xxxxxxxxxxxxxxxxxxxxxxxxxxxxx\=\kill
\> \textbf{repeat} $2n$ times:\\
\> 1.\> \textbf{if} $v$ odd \textbf{and} ($u$ even \textbf{or} $u>v$): \> swap $(u,v)$, swap $(r,s)$; remember it\\
\> 2.\> \textbf{if} $u,v$ both odd: \> $v \mathrel{-}= u$; \quad $s \mathrel{+}= r$\\
\> 3.\> \> $v \mathrel{/}= 2$; \quad $r \mathrel{\times}= 2$ \ (the doubling only if the round fired)\\
\> 4.\> \textbf{if} it swapped: \> swap back
\end{tabbing}
\end{minipage}
\end{center}

\noindent
Started at $(u,v,r,s)=(p,x,0,1)$ this terminates with $u=1$, $v=0$, $s=p$ and
\begin{equation}\label{ec:eq:pseudo}
r \;=\; -x^{-1}\,2^{k} \bmod p,
\end{equation}
where $k$ counts the rounds that did anything. That is a \emph{pseudo}-inverse: right up to a sign
and a power of two, both of which the caller fixes --- one modular negation, and $k$ modular
halvings driven by a counter. Equation~\eqref{ec:eq:pseudo} is asserted directly in
\code{tests/test\_ec\_kaliski.py} for every invertible $x$ and every $p$ tested, alongside
$u=1$, $v=0$, $s\equiv0$ and the counter itself.

\begin{keybox}[title={Three bits per round, and why none can be dropped}]
The round is a permutation of $(u,v,r,s)$ chosen by its decision bits, so those bits must survive
or the round is not reversible. Three are kept:
\begin{itemize}\setlength\itemsep{1pt}
\item $f$ --- \emph{did this round do anything} ($v\neq0$ on entry). Not recoverable afterwards: an
inert round and the last active round both leave $v=0$.
\item $a$ --- \emph{did it swap}.
\item $b$ --- \emph{were both odd}, so the subtract-and-add fired. Not implied by $(f,a)$: the case
$f=1,a=0$ covers both ``$v$ was even'' and ``both odd with $u\le v$''.
\end{itemize}
The comparison $u>v$ is \emph{not} kept. It is consumed while building $a$ and then uncomputed by
running the comparator a second time --- before $u$ and $v$ change, which is the only window in
which that is possible.
\end{keybox}

\noindent
In code, with the decision bits built first and the comparison uncomputed before $u$ and $v$ move:

%%LST ec_kaliski.kaliski_round drop elide=t,~sc~=~m.anc(1,~"t")..m.free(t,~sc)%%

\noindent
$r$ and $s$ are reduced mod $p$ every round. They have to be: $r$ doubles each round, and the tests
confirm it runs past $p$ within a few rounds if left alone. $u$ and $v$ are plain $n$-bit integers,
staying in $[0,p]$ throughout.

\paragraph{Division, and where the garbage goes.}
What a point addition actually wants is $\lambda = y/x$, and the inversion above leaves $6n$ qubits
of per-round records behind. Both are solved at once by the shape of \cite{ec:rnsl17}: invert,
multiply, copy the answer out, then \emph{run the whole thing backwards}. Only the copy survives;
records, counter, $u$, $v$ and $s$ all return to $\ket0$, which the simulator checks at the exact
instruction each is freed.

%%LST ec_kaliski.mod_div doc=1%%

\begin{warnbox}[title={A control that has to reach all the way down --- here, though not everywhere}]
The division is controlled by the counting qubit $q$, and for the inversion built above it is not
enough to gate only the final copy. On the $q=0$ branch the $x$ register does not hold $x_1-x_2$;
by that point in the point addition it holds $x_1+2x_2$, which vanishes for perfectly ordinary
accumulators, so the circuit would be inverting zero --- and Kaliski on $0$ terminates with $u=p$
and $s=1$, not the $u=1$, $s\equiv0$ that the cleanup here relies on to return those registers to
$\ket0$.

The fix costs almost nothing. Load the inversion's \emph{setup} under the control ---
$u\leftarrow p$, $v\leftarrow x$, $s\leftarrow1$ --- and when $q=0$ every register starts at zero,
so $v=0$ makes all $2n$ rounds inert, $r$ stays $0$, and the correction does nothing. The rounds
themselves are never controlled; $n$ ANDs and a handful of CNOTs buy the whole thing, and it is
checked directly: with $q=0$ the output register is $0$ for every input.

\cite{ec:kim26} does \emph{not} do this, and does not need to. Its Fig.~7 runs the inversion and
the multiplication unconditionally and gates only the copy-out of $\lambda$, which is sound there
because the inversion is undone by its literal inverse --- whatever Inv did to the registers on a
degenerate input, Inv$^\dagger$ undoes. The gating above is the price of \emph{this} construction's
cheaper cleanup, not a correction to the paper.
\end{warnbox}

\section{Level 4: the eleven-step point addition}\label{ec:sec:padd}

The accumulator $(x_1,y_1)$ is quantum, the addend $(x_2,y_2)$ classical, and the control $q$ says
whether to add at all. In place, no garbage:

\begin{center}\small\setlength{\tabcolsep}{6pt}
\begin{tabular}{@{}rlll@{}}
\toprule
& operation & \multicolumn{2}{l}{after it, on the $q=1$ branch:}\\
\cmidrule(l){3-4}
& & $x$ register & $y$ register\\
\midrule
1 & $x \mathrel{-}= x_2$                      & $x_1-x_2$ & $y_1$\\
2 & $y \mathrel{-}= q\,y_2$                   & & $y_1-y_2$\\
3 & $\lambda \mathrel{\oplus}= q\cdot y/x$    & & \\
4 & $y \mathrel{\oplus}= x\lambda$            & & $\mathbf{0}$\\
5 & $x \mathrel{+}= 3x_2$                     & $x_1+2x_2$ & \\
6 & $x \mathrel{-}= \lambda^2$                & $x_2-x_3$ & \\
7 & $y \mathrel{+}= x\lambda$                 & & $\lambda(x_2-x_3)$\\
8 & $\lambda \mathrel{\oplus}= q\cdot y/x$    & & \\
9 & $x \mathrel{-}= (2-q)x_2$                 & $-x_3$ & \\
10 & $y \mathrel{-}= q\,y_2$                  & & $y_3$\\
11 & $x \leftarrow -x$ \ if $q$               & $x_3$ & \\
\bottomrule
\end{tabular}
\end{center}

\noindent
The circuit is the table, line for line:

%%LST ec_pointadd.point_add_ctrl doc=1%%

\noindent
Three things in that table are worth pausing on.

\emph{Step 4 is an XOR, not a subtraction.} At that moment $y$ holds $y_1-y_2$ and $x\lambda$ is by
definition $y_1-y_2$, so XOR-ing the product in zeroes the register, and the algebra guarantees it.
What is saved is the \emph{accumulation} --- CNOTs instead of a modular subtraction. The product
$x\lambda$ still has to be computed and uncomputed around it, which is two modular multiplications:
the step is cheaper than the alternative, not cheap.

\emph{Step 8 is the same circuit as step 3, and it \emph{clears} $\lambda$.} By then the registers
hold $x=x_2-x_3$ and $y=\lambda(x_2-x_3)$, so $y/x$ is again $\lambda$, and the XOR-accumulating
division undoes what step 3 wrote. The whole eleven-step arrangement exists to make that identity
true; without it $\lambda$ would be $n$ qubits of garbage per addition.

\emph{Step 9 is two of \cite{ec:hjnrs20}'s steps merged}, following \cite{ec:kim26}. The original
does $x\mathrel{-}=2x_2$ then a controlled $x\mathrel{+}=x_2$; the merged form subtracts
$(2-q)x_2$ once. Building that operand is free: $2x_2\bmod p$ and $x_2\bmod p$ are both classical,
so the register is X-gated to one of them and CNOTs from $q$ patch the bits where they differ.

\begin{warnbox}[title={Three exceptional cases, not one}]
The circuit is correct except on a set of density $O(1/p)$, and it does not detect the exceptions
--- a check would cost more than it saves. There are three, and only the first is the one usually
quoted:
\begin{itemize}\setlength\itemsep{1pt}
\item $x_1=x_2$ --- step 3 divides by $x_1-x_2$. Doubling, adding an inverse, or infinity.
\item $x_3=x_2$, i.e.\ $P_1=-[2]P_2$ --- step 8 divides by $x_2-x_3$. This one is invisible in
\eqref{ec:eq:add} and appears only once you follow what the \emph{registers} hold. The in-place
multiplier of \S\ref{ec:sec:dialog} reaches the same wall from the other direction, where it would
be asked to multiply by zero.
\item for the \emph{windowed} circuit of \S\ref{ec:sec:window} adding $\mathcal O$: $x_R=0$. This
one is not about the table above --- it belongs to \cite{ec:schrott26}'s Algorithm~1, whose own
step~6 divides by $x_R-x_{P_i}$, and that is $x_R$ when $P_i=\mathcal O$ is encoded as $(0,0)$.
Adding the identity looks like it could never fail.
\end{itemize}
\code{ec\_classical.point\_add\_exceptional} decides the first two; the tests exclude them
explicitly and report how many, rather than passing quietly on inputs they never exercise.
\end{warnbox}

\section{Making it fast}\label{ec:sec:fast}

Everything from here on is measured. Each subsection fixes a \emph{task} and varies one
implementation choice; comparing a modular doubling against a point addition would say nothing, so
nothing here does. Every variant is re-checked against the classical model \emph{inside} the
benchmark before its cost is recorded --- a cheaper circuit that is wrong is not a datapoint ---
and variants that are deliberately approximate record a measured failure rate instead. All of it
comes from \code{bench/ec\_ablation.py}, and the tables below are generated from its output rather
than typed.

\begin{keybox}[title={Which Toffoli count, and why it matters here more than in Part V}]
Both source papers count ``Toffolis'' as CCX, CCZ and AND gates together
\cite{ec:schrott26}. A measurement-based AND$^\dagger$ contains \emph{no} Toffoli and no $T$ gate,
so it is not counted. That is the convention used in every table below, and it is the one to
compare against published figures.

It is not the only defensible convention --- one could count every three-qubit operation performed,
uncomputes included --- and the two differ by nearly $2\times$ on Gidney-style circuits, which is
most of what this part builds. \code{ec\_cost} reports both (\code{toffoli\_paper} and
\code{toffoli\_equiv}) precisely so the choice is visible rather than buried. Quoting the wrong one
would make every optimisation here look roughly half as good as the literature says it is.
\end{keybox}

\subsection{Level 1: which adder, and how to control it}\label{ec:sec:fastadd}

Two adders, chosen per call site by how many ancillas happen to be free --- as \cite{ec:schrott26}
puts it, the ancilla-hungry one where ancillas are available and the ancilla-free one where they
are not.

\begin{center}\small
\begin{tabular}{@{}rrrrrrrr@{}}
\toprule
& \multicolumn{2}{c}{addition} & \multicolumn{3}{c}{controlled addition} & \multicolumn{2}{c}{qubits}\\
\cmidrule(lr){2-3}\cmidrule(lr){4-6}\cmidrule(lr){7-8}
& & & & \multicolumn{2}{c}{copy-then-add, with} & & \\
\cmidrule(lr){5-6}
$n$ & CDKM & Gidney & reconstruct & CDKM & Gidney & CDKM & Gidney\\
\midrule
%%ADDERS%%
\bottomrule
\end{tabular}
\end{center}

\noindent
Plain addition is the $2n$-versus-$n$ of the textbook comparison, and lands exactly there:
$@@cdkm32@@=2n$ against $@@gid32@@=n-1$ at $n=32$. CDKM spends two Toffolis per bit, Gidney one AND
per bit, and the difference is that Gidney's targets start clean, so they uncompute for free --- the
$n-1$ rather than $n$ being the carry into the top bit, which no adder has to compute.

Controlled addition is where \cite{ec:kim26} makes its choice. Promoting the adder's CNOTs to
Toffolis costs $@@recon32@@$ at $n=32$; copying the operand into clean ancillas \emph{under the
control} and then adding unconditionally costs $@@copy32@@$, a factor $@@cadd_ratio@@$. The copies
all commute, so they also parallelise, which the reconstructed version cannot.

%%LST ec_adders.cadd doc=1%%

\noindent
The middle column is there to keep that factor honest. Copy-then-add and the Gidney adder are two
separate decisions, and the $@@cadd_ratio@@\times$ bundles them; hold the adder fixed at CDKM and
the control style alone is worth $@@cadd_fixed@@\times$ ($@@copyc32@@$ against $@@recon32@@$).
The rest is the adder.

What it costs is qubits. The last two columns are the plain adders' widths, $2n+1$ against
$3n-1$; for the controlled pair the gap is wider still, $@@recon32q@@$ against $@@copy32q@@$ at
$n=32$, since the copy needs a register of its own.

\subsection{Level 0: temporary ANDs}\label{ec:sec:fastand}

Part~V measured this for factoring; it applies unchanged, and matters more here because there is so
much more arithmetic to apply it to. Over one whole modular inversion:

\begin{center}\small
\begin{tabular}{@{}rrrrrr@{}}
\toprule
$n$ & ANDs & AND$^\dagger$s & $T$ (all Toffoli) & $T$ (temporary AND) & saving\\
\midrule
%%TEMPAND%%
\bottomrule
\end{tabular}
\end{center}

\noindent
A flat $@@ta_ratio@@\times$, and the arithmetic says why: ANDs and AND$^\dagger$s arrive in
near-equal numbers, so the average is $(4+0)/2=2$ against $7$, i.e.\ $3.5\times$ on the ANDs alone,
diluted by the Fredkins and dirty-target Toffolis that cannot use the trick.

\subsection{Level 2: approximate and pseudo-Mersenne reduction}\label{ec:sec:approx}

Shor's circuit only has to work with large constant probability, and \cite{ec:schrott26} spends
that freedom on the reductions. Two ideas.

\emph{Compare only the top bits.} ``Is this at least $p$?'' is answered correctly unless the top
few bits tie, which happens with probability about $2^{-\mathrm{msbs}}$. That shrinks an $n$-bit
comparison to an $\mathrm{msbs}$-bit one and --- more importantly --- lets the flag be produced
\emph{before} the subtraction rather than recovered after it, collapsing the exact two-step
reduction of \S\ref{ec:sec:modarith} into a single controlled subtraction.

\emph{Exploit the shape of the prime.} secp256k1 uses $p=2^{u}-f$ with $f=2^{32}+977$, tiny against
$2^{256}$. Then ``at least $p$?'' is almost ``is bit $u$ set?'', and subtracting $p$ becomes ``clear
bit $u$, then add $f$'' --- and since $f$ is small, adding it need only touch the bottom few bits,
because the carry cannot travel far. An $n$-bit constant subtraction becomes a short one.

The pseudo-Mersenne doubling is the whole idea in eight lines --- note that nothing compares
anything, and that clearing the overflow bit \emph{is} the subtraction of $2^u$:

%%LST ec_approx.moddbl_pm drop%%

\begin{center}\small
\begin{tabular}{@{}rlrrrr@{}}
\toprule
& & \multicolumn{2}{c}{doubling} & \multicolumn{2}{c}{controlled addition}\\
\cmidrule(lr){3-4}\cmidrule(lr){5-6}
$p$ & reduction & Toffoli & failure & Toffoli & failure\\
\midrule
%%MODARITH%%
\bottomrule
\end{tabular}
\end{center}

\noindent
Read the failure column as a rate, not a defect: these circuits are \emph{supposed} to be wrong
sometimes. The question is whether the rate tracks $2^{-\mathrm{msbs}}$, and it does --- at
$p=4093$, four MSBs give $@@m4@@\%$ and eight give $@@m8@@\%$, a factor $@@m_ratio@@$ for four
bits. The pseudo-Mersenne doubling is the headline: $@@dbl_pm@@$ Toffolis against $@@dbl_exact@@$
exact, a $@@dbl_pm_ratio@@\times$ reduction, for a $@@dbl_pm_fail@@\%$ failure rate at this width.

\begin{warnbox}[title={Failure compounds, which is why the paper's budget is so large}]
A single operation failing a few per cent of the time sounds cheap. Chained through the
$\approx1.4n$ iterations of an in-place multiplication it is not. Measured end to end at $p=31$: a
two-bit comparison, which fails about $12\%$ per operation, turns a working multiplier into one that
fails $81\%$ of the time, while the full-width comparison stays exact. That is the calculation behind \cite{ec:schrott26}'s choice of $40$ to $50$
bits at $n=256$ --- per-operation failure near $2^{-40}$, so even several hundred iterations stay
negligible. A cheap approximation is only cheap once you count how often you use it. The measurement
is in \code{tests/test\_ec\_opt1128.py}, not an estimate.
\end{warnbox}

\subsection{Level 3: Montgomery multiplication over a carry-save tree}\label{ec:sec:mont}

\cite{ec:kim26} rewrites word-level Montgomery so that its two halves have the same shape. The
textbook reduction computes $M=c_w p'_w \bmod 2^w$ and then $c \mathrel{+}= Mp$; substituting $M$
gives $c \mathrel{+}= c_w\,(p'_w \bmod 2^w)\,p$, where the bracket depends only on the modulus. So
$d=(p'_w\bmod 2^w)\,p$ is a compile-time constant and the reduction becomes a second
``$c \mathrel{+}= c_w\cdot d$'', structurally identical to the multiplication step. One carry-save
tree then serves both. Its 3:2 compressor turns three addends into two with no carry propagation at
all,
\[
\begin{aligned}
S_i &= A_i \oplus B_i \oplus C_i && (\text{three CNOTs}),\\
M_i &= \MAJ(A_i,B_i,C_i) = A_i \oplus \big((A_i{\oplus}B_i)\wedge(A_i{\oplus}C_i)\big) && (\text{one AND}),
\end{aligned}
\]
with $A+B+C=S+2M$, so reducing $w$ addends to two costs $O(\log w)$ depth and only the final
addition has to propagate a carry.

%%LST ec_qcsa.csa_compress doc=1%%

\begin{center}\small
\begin{tabular}{@{}lrrrr@{}}
\toprule
multiplier & Toffoli & qubits & depth & depth vs.\ schoolbook\\
\midrule
%%MUL%%
\bottomrule
\end{tabular}
\end{center}

\begin{warnbox}[title={What this table does \emph{not} show, and why}]
The carry-save multiplier is a \emph{depth} optimisation, and depth is the column where it wins:
$@@mul_depth_qcsa@@\times$ over schoolbook at $n=16$. It does not beat the QROM-lookup multiplier
of \cite{ec:hjnrs20} on gate count, and at these widths it should not be expected to --- a tree
over two or four words has almost nothing to flatten, and the rewrite costs real width. The two differ by
$2^w\lfloor c_wp'_w/2^w\rfloor\,p$, which with $c_w$ and $p'_w$ both below $2^w$ reaches about
$2^{2w}p$ --- so the reduction over-adds by that much and the accumulator settles at $n+2w$ bits
rather than the $n+2$ of textbook Montgomery, with a final reduction to mop up. The paper does not dwell
on this; the tests measure the width rather than trusting it.
\end{warnbox}

\subsection{Level 3: dropping the counter and postponing the reduction}\label{ec:sec:kalopt}

Two changes to \S\ref{ec:sec:kaliski}, both from \cite{ec:kim26}, and they compound.

\emph{Run every round unconditionally.} Drop the ``is $v$ still nonzero?'' test. An inert round does
nothing except double $r$, so instead of stopping at $-x^{-1}2^{k}$ the register runs on to
$-x^{-1}2^{2n}$ --- the same value for \emph{every} input, with no counter to keep and no $2n-k$
corrective halvings to perform. And $2^{2n}$ is exactly right: fed $x=X2^{n}$ in Montgomery form the
answer is $-X^{-1}2^{n}$, the Montgomery form of $-X^{-1}$, so the power of two has corrected itself
and only the sign is left --- one modular negation, against \cite{ec:hjnrs20}'s counter and its
$2n-k$ controlled halvings. It also removes one of the three recorded bits per round, since $f$ existed only to
say whether the round fired.

\emph{Postpone the modular reductions.} Do not reduce $r$ and $s$ mod $p$ each round; let them run
into a $2n$-bit register and reduce once at the end. This changes what a round \emph{is}: a modular
doubling (compare, conditionally subtract, clear the flag) becomes a plain left shift, which is a
renaming of wires and costs no Toffoli at all, and a modular addition becomes a plain addition.

\begin{center}\small
\begin{tabular}{@{}rrrrrrrrr@{}}
\toprule
& \multicolumn{3}{c}{Toffoli} & \multicolumn{3}{c}{depth} & \multicolumn{2}{c}{qubits}\\
\cmidrule(lr){2-4}\cmidrule(lr){5-7}\cmidrule(lr){8-9}
$n$ & \cite{ec:hjnrs20} & \cite{ec:kim26} & saved & \cite{ec:hjnrs20} & \cite{ec:kim26} & saved & \cite{ec:hjnrs20} & \cite{ec:kim26}\\
\midrule
%%INV%%
\bottomrule
\end{tabular}
\end{center}

\noindent
At $n=12$: $@@inv_ref@@$ Toffolis become $@@inv_opt@@$, a $@@inv_toff_pct@@$ saving, and depth falls
by $@@inv_depth_pct@@$, for about $11\%$ more qubits.

\cite{ec:kim26} reports $58$--$60\%$ for division at $n=192$--$521$, but on a different metric:
the \emph{product} of qubit count and $T$-depth, which is the trade-off measure that paper uses
throughout. The columns above are plain counts, which is why they are not directly comparable ---
a construction that spends qubits to buy depth looks better under a product metric than under
either factor alone. The direction agrees; the magnitudes should not be read against each other.

The $2n$-bit budget is not a comfortable margin, it is the exact requirement --- $r$ reaches $6$
bits for $p=7$ and $32$ bits for $p=65521$, both exactly $2n$. The test measures it rather than
assuming it.

\subsection{Level 3, rebuilt: the Euclidean dialog}\label{ec:sec:dialog}

Everything so far treats inversion and multiplication as two jobs. Schrottenloher
\cite{ec:schrott26}, following \cite{ec:khattar25}, observes that they are one, and this is the
single largest structural saving in this part.

Run the binary GCD on $(u,v)$ alone and \emph{do not} maintain the B\'ezout coefficients. Just
record the two decision bits per iteration --- call the record a \emph{dialog}. It is $1.5$ bits
per iteration on average, against the $2n$ qubits carrying $(r,s)$ would cost:

\begin{center}\small
\begin{minipage}{0.82\textwidth}
\begin{tabbing}
xxx\=xxxx\=xxxxxxxxxxxxxxxxxxxxxxxxxxxxxxxxxxxx\=\kill
\> \textbf{repeat} $\approx1.413n$ times:\\
\> 1.\> $b_0 := v \bmod 2$; \quad $b_1 := [\,u>v\,]$ \> record $b_0$ and $b_0b_1$\\
\> 2.\> \textbf{if} $b_0b_1$: \ swap $(u,v)$\\
\> 3.\> \textbf{if} $b_0$: \ $v \mathrel{-}= u$\\
\> 4.\> $v \mathrel{/}= 2$
\end{tabbing}
\end{minipage}
\end{center}

\noindent
The iteration count is not arbitrary. Halving $v$ always removes a bit; the subtraction removes a
second one half the time. \cite{ec:schrott26} averages the shrink factor of $uv$ over those two
cases, $\tfrac12\cdot\tfrac12+\tfrac12\cdot\tfrac14=\tfrac38$, and reads off
$2n/\log_2(8/3)\approx1.413n$ iterations. That step is heuristic --- it averages the factor rather
than its logarithm, which is not the same thing --- but the conclusion holds: measured on this
implementation's own dialog, the mean iteration count is $1.383n$ at $n=32$ and $1.409n$ at
$n=256$, climbing towards $1.413n$. The spread is $\approx0.6\sqrt n$, which is what the $+c\sqrt n$
safety margin is for.

%%LST ec_eea.dialog_round doc=1%%

\noindent
Now the second idea. The updates to $(r,s)$ that the full extended algorithm would have performed
are all \emph{linear} and depend only on those two bits, so they can be replayed afterwards, onto
any starting pair:
\[
s \leftarrow 2s \bmod p; \qquad \text{if } b_0:\ s \mathrel{+}= r; \qquad \text{if } b_0b_1:\ \text{swap}(r,s)
\]
Read \emph{forwards}, alongside the Euclidean algorithm as the textbook extended version would,
this maps $(1,0)\mapsto(0,\,x^{-1})$ --- an inversion --- and more generally
$(y,0)\mapsto(0,\,yx^{-1})$, which is an inversion and a multiplication for the price of one pass.
That is \cite{ec:schrott26}'s observation, and it is already the saving.

But the circuit reads the record \emph{backwards}, and gets a different map:
$(y,0)\mapsto(0,\,yx)$. Multiplication by $x$, not by $x^{-1}$. That is not a defect but the
reason the backwards reading is chosen: replayed in reverse the updates need only \emph{doubling}
and addition, where the forward reading needs modular \emph{halving}. And since the whole thing is
a circuit, running it in the other direction recovers the division. One component, both of the
point addition's multiplicative steps.

The replay is seven lines, and the whole saving is that none of them looks at $x$:

%%LST ec_eea.bezout_replay doc=1%%

\noindent
Assembling the three phases --- dialog, replay, dialog backwards --- gives the in-place
multiplication, and the rebuild consumes the record on its way out, so nothing is left behind:

%%LST ec_eea.inplace_mul doc=2%%

\begin{keybox}[title={Direction matters, and it is easy to get wrong}]
Applying Algorithm~3's update rules in the record's original order is \emph{not} the inverse of
applying them in reverse, and it is not the transpose either --- it is simply a different linear
map. Writing ``$(1,0)$ gives $x^{-1}$'' next to ``$(y,0)$ gives $yx$'' is self-contradictory, since
the second is linear in $y$ and contains the first at $y=1$; an earlier draft of this section said
exactly that, and the code is what caught it.

The classical model in \code{ec\_classical.py} carries all three maps and the tests pin down which
is which: \code{bezout\_replay} (reversed) multiplies by $x$, \code{bezout\_replay\_inv} divides by
it, and \code{bezout\_inverse\_classical} --- an independent coefficient-tracking implementation
that needs modular halving where the replay needs only doubling --- yields $x^{-1}$ directly. The
third exists purely as a cross-check, which is why the direction claim is not taken on faith.
\end{keybox}

\noindent
Compared against what it replaces --- Kaliski inversion plus an out-of-place multiplication, which
is what a point addition otherwise needs:

\begin{center}\small
\begin{tabular}{@{}rrrrrrr@{}}
\toprule
& \multicolumn{2}{c}{invert-then-multiply} & \multicolumn{2}{c}{dialog + replay} & \multicolumn{2}{c}{saving}\\
\cmidrule(lr){2-3}\cmidrule(lr){4-5}\cmidrule(lr){6-7}
$p$ & Toffoli & qubits & Toffoli & qubits & Toffoli & qubits\\
\midrule
%%DIALOG%%
\bottomrule
\end{tabular}
\end{center}

\noindent
At $p=61$ that is $@@dlg_toff_pct@@$ of the gates and $@@dlg_q_pct@@$ of the qubits, and the
margin widens with $n$. The record itself is
squeezed further by a trick worth its own paragraph: each pair $(b_0, b_0b_1)$ can only be
$(0,0)$, $(1,0)$ or $(1,1)$ --- never $(0,1)$, since $b_0b_1$ cannot be set unless $b_0$ is --- so
three pairs carry $3^3=27<32$ states and fit in five bits instead of six. One qubit freed per three
iterations, taking the record from $2.83n$ to $2.355n$ asymptotically, and (unlike a shift-based
encoding) leaving it a \emph{fixed} size, which is what lets the rest of the circuit be laid out
statically.

\subsection{Level 4, rebuilt: projective coordinates}\label{ec:sec:proj}

The other way to make an inversion cheap is not to do it. In Jacobian coordinates
$(X:Y:Z)\sim(x,y)=(X/Z^2,\,Y/Z^3)$ the addition law is inversion-free, and with the addend in
affine form ($Z_2=1$, which it is, being a classical point) every term involving $Z_2$ drops out:
\[
\begin{aligned}
U_2 &= x_2Z_1^2, & S_2 &= y_2Z_1^3, & H &= U_2-X_1, & R &= S_2-Y_1,\\
X_3 &= R^2-H^3-2X_1H^2, & Y_3 &= R(X_1H^2-X_3)-Y_1H^3, & Z_3 &= Z_1H. &&
\end{aligned}
\]
Eleven multiplications, no division --- \cite{ec:kim26} reports this as the fewest of any quantum
projective addition, against a previous best of twelve.

The catch is that projective coordinates are not a unique representation. $(P+Q)-Q$ returns a
\emph{different triple} for the same point, so the input cannot be erased and every addition leaves
its accumulator behind. The fifteen intermediates \emph{can} all be cleaned --- each by reversing
the step that made it, with that step's inputs still live --- and this implementation cleans them,
but the input point is $3n$ qubits of irreducible garbage per addition.

The eighteen steps, and then the unwinding --- note the last two lines, which are the whole
subtlety:

%%LST ec_proj.jacobian_add doc=1 elide=def~mul(a,~b,~out)..MA.modsub(m,~b,~out,~p) elide=Hsq~=~A(n,~"Hsq")..twoV~=~A(n,~"2V")%%

\begin{warnbox}[title={The order of unwinding is not simply ``backwards''}]
$X_3$ is an output and must survive, so the intermediate $T_0$ cannot be released by reversing the
step that consumed it. It is released by reversing the step that \emph{produced} it, from $R^2$ and
$H^3$, which are still live at that moment. Getting this wrong destroys the answer;
\code{Machine.undo} exists precisely so that a recorded step can be reversed out of order, and the
tests sweep the projective representatives $Z$ of every input point --- $Z\in\{1,2\}$ by default
and all of $[1,p)$ under \code{SHOR\_EC\_FULL=1} --- a sweep that catches representation bugs an
affine-only test cannot see.
\end{warnbox}

\subsection{Level 5: windowing, zig-zag, and the semiclassical transform}\label{ec:sec:window}

Three counting arguments, none of which changes a single arithmetic circuit.

\emph{Windowing} \cite{ec:gidney-ekera} groups $w$ control qubits into one address and looks up
$P_i$ from a table of $2^w$ precomputed points, so $2n+2$ additions become $2\lceil (n{+}1)/w\rceil$.
A lookup costs $2^w$ Toffolis, an order of magnitude below an addition, so $w$ is chosen large ---
both \cite{ec:schrott26} and \cite{ec:babbush26} use $w=16$.

\emph{Zig-zag} \cite{ec:grassl16} attacks the projective garbage. Write $\rho$ for the number of
accumulator registers --- not $m$, which above is the counting-register width. Additions are
arranged in runs of $\rho, \rho-1, \dots, 1$: fill all $\rho$ registers, then walk back
un-computing all but the last of the run --- legal, because each state's predecessor is still live
--- returning $\rho-1$ registers, and repeat. Total forward additions $\rho(\rho+1)/2$, so $\rho$
registers serve them all.

\emph{The semiclassical transform} \cite{ec:griffiths-niu} is the same trick as
\S\ref{sec:onectrl}: every controlled rotation in the inverse QFT is controlled by a qubit about to
be measured, so measure it first and make the control classical. Both scalar registers collapse to
\emph{one recycled qubit}. This is the assumption behind both papers reporting only the cost of the
point arithmetic.

The windowed addition itself, with the three lookups and the two in-place multiplications that
replace \S\ref{ec:sec:padd}'s two divisions:

%%LST ec_window.windowed_point_add doc=1 elide=xs~=~[0..x3s~=~[0 tail=MA.cmodneg(m,~c[0],~x2,~p)%%

\begin{center}\small
\begin{tabular}{@{}rrrrrrrrr@{}}
\toprule
& \multicolumn{3}{c}{point additions} & \multicolumn{3}{c}{zig-zag registers} & \multicolumn{2}{c}{garbage qubits}\\
\cmidrule(lr){2-4}\cmidrule(lr){5-7}\cmidrule(lr){8-9}
$n$ & none & $w{=}11$ & $w{=}16$ & none & $w{=}11$ & $w{=}16$ & none & $w{=}16$\\
\midrule
%%WINDOW%%
\bottomrule
\end{tabular}
\end{center}

\noindent
At $n=256$: $@@w256_none@@$ additions become $@@w256_16@@$, a $@@w256_ratio@@\times$ reduction,
and the registers holding projective garbage fall from $@@zz256_none@@$ to $@@zz256_16@@$. The two
compose --- windowing cuts how many additions there are, zig-zag cuts how many of their outputs are
live at once --- and together they take the garbage from $@@g256_none@@$ qubits to $@@g256_16@@$.

\section{Everything at once}\label{ec:sec:cost}

The subsections above each hold one thing fixed. This one does not: it stacks them, on the operation
that actually matters --- one point addition --- and then on the whole circuit.

\begin{center}\small
\begin{tabular}{@{}llrrrr@{}}
\toprule
curve & construction & Toffoli & qubits & depth & vs.\ affine\\
\midrule
\multicolumn{6}{@{}l}{\footnotesize Alg.~3, Alg.~4 are \cite{ec:kim26}'s; Alg.~1 is \cite{ec:schrott26}'s.}\\
%%PADD%%
\bottomrule
\end{tabular}
\end{center}

\noindent
Three constructions of the same map, and they stack differently. The affine addition is in place and
leaves nothing behind, and pays for that with an inversion. The Jacobian addition removes the
inversion entirely and is $@@padd_proj_pct@@$ cheaper, but leaves $3n$ qubits of input point as
garbage --- which is only affordable once windowing and zig-zag have cut how many such inputs are
live at once. The windowed addition is cheapest of the three at $@@padd_win_pct@@$ below affine,
and it gets there differently again: not by changing coordinates, but by replacing
invert-then-multiply with a single Euclidean dialog, and by amortising three table lookups over a
whole window of additions.

\begin{keybox}[title={How the savings compose --- and where they do not}]
Read down the stack for one point addition at $n=256$, taking each factor from the measurements
above:
\begin{itemize}\setlength\itemsep{2pt}
\item \emph{Gate-level, and only partly independent.} The Gidney adder ($\approx2\times$ on the
inner additions) and the temporary AND ($@@ta_ratio@@\times$ on $T$) are \emph{not} two savings:
they are one saving counted at two levels. The adder wins by choosing a gate whose target starts
clean, and the temporary-AND accounting is exactly what makes that gate's uncompute free.
Multiplying the two would double-count. Copy-then-add controls are a smaller further factor than
the table suggests: the $@@cadd_ratio@@\times$ there bundles the control style together with the
switch to Gidney's adder inside it. Hold the adder fixed and the control style alone is worth
$@@cadd_fixed@@\times$ ($@@copyc32@@$ against $@@recon32@@$ at $n=32$, CDKM on both sides) --- it
buys parallelism and depth more than gates. Pseudo-Mersenne reduction
($@@dbl_pm_ratio@@\times$ on doublings) is the one that cleanly composes, because it acts on a
different part of the circuit again: the reductions rather than the additions.
\item \emph{Structural, and they replace rather than compose.} The dialog
($@@dlg_toff_pct@@$ off the multiplicative step) and projective coordinates
($@@padd_proj_pct@@$ off the addition) are two different answers to the same problem --- the
inversion --- so you take one, not both. Every published construction picks a side.
\item \emph{Count-level, and these do multiply the whole thing.} Windowing divides the number of
point additions by $@@w256_ratio@@$ at $n=256$; zig-zag divides the live projective garbage by
$@@zz_ratio@@$; the semiclassical transform deletes both counting registers. None of them touches an
arithmetic circuit, which is precisely why they compose cleanly with everything above.
\end{itemize}
The last group dominates. An optimisation that halves a point addition is worth half as much as one
that runs fifteen times fewer of them, which is why both papers spend their first section on
windowing and their last on gate counts.
\end{keybox}

\noindent
And the whole circuit, on the toy curve $y^2=x^3+5x+4 \bmod 7$ with $\mathrm{ord}(G)=5$:

\begin{center}\small
\begin{tabular}{@{}lrrrr@{}}
\toprule
ECDLP circuit & qubits & Toffoli & gates & depth\\
\midrule
%%FULL%%
\bottomrule
\end{tabular}
\end{center}

\noindent
$@@full_q@@$ qubits and $@@full_g@@$ gates to compute a discrete logarithm in a group of order
five. That ratio is the honest summary of this part: the arithmetic is enormous, and the
counting-level tricks are the only reason cryptographic sizes are discussed at all. Note the last
two rows --- the semiclassical transform takes $@@tab_q@@$ qubits to $@@sc_q@@$ by recycling one
control qubit, which is precisely why both papers report only the point arithmetic and treat the
scalar registers as free.

\subsection{Space at cryptographic size}\label{ec:sec:space}

The one place where a projection is worth making, because it is what the papers' headline qubit
counts are built from. The dialog record dominates the space of \cite{ec:schrott26}'s construction:

\begin{center}\small
\begin{tabular}{@{}rrrrr@{}}
\toprule
$n$ & record, raw & with Fig.~1 packing & per $n$ & with register sharing\\
\midrule
%%SPACE%%
\bottomrule
\end{tabular}
\end{center}

\noindent
The middle column converges to the paper's $2.355n$ from above --- the gap is the $O(\sqrt n)$
safety margin on the iteration count, which is a constant fraction at these sizes and vanishes
asymptotically. The last column is \emph{not} implemented here: it is Sec 3.1's register sharing, in
which $u$ and $v$ shrink as the algorithm proceeds and the freed qubits absorb the record. It is
priced here but not built --- \S\ref{ec:sec:gaps} says why.

The paper's own closed form for the total --- its asymptotic $2.355n$ record plus $2n$ for the pair
of modular integers the reconstruction holds, i.e.\ $4.355n+O(\sqrt n)$ --- gives
$@@q1128_space@@$ qubits at $n=256$ against the paper's $1192$, and $@@q1128_gate@@$ against $1446$
for the gate-optimised variant. That is the paper's formula evaluated, \emph{not} the table above
summed: the table's middle column carries its $O(\sqrt n)$ margin explicitly and would give a
larger figure. Neither is a check on the other.

The full-algorithm count of $2^{@@shor_log2@@}$ Toffolis is likewise the paper's equation~(1)
evaluated at the paper's own per-addition figure --- restated for scale, not independently
confirmed. Nothing in this package measures a $256$-bit point addition.

\section{Verifying it}\label{ec:sec:verify}

A wrong ECDLP circuit does not return a wrong answer. It returns a \emph{flat histogram} --- and
since the classical post-processing verifies its own candidates against the curve, it will happily
print the correct $k$ from pure noise. Testing the output is therefore worth almost nothing; three
other things are worth a great deal.

\paragraph{Exact simulation at full width.}
Every arithmetic circuit here is a permutation of basis states, so pushing one basis state through
in $O(\text{gates})$ time is exact --- and scales to hundreds of qubits where a statevector dies at
thirty. Ancillas are checked back to $\ket0$ \emph{at the instruction where they are freed}, not
merely at the end, which is what catches a circuit that computes the right value and leaves a dirty
flag behind. This is sound only for circuits whose phases carry no information; the sole phase and
Hadamard gates outside the AND gadget are in the state preparation and the Fourier transform, which
are run on Aer instead.

%%LST ec_sim.simulate doc=2 elide=by_pos~=~{}..def~apply(op,~qargs) tail=bits[w[2]]~=~0%%

\paragraph{Statevector cross-checks.}
The fast simulator only sees basis states, and a bug in it would hide from every test that uses it.
So the small circuits are also evolved on genuine superpositions and checked for exact fidelity
\emph{and} for the absence of leaked amplitude --- the statement that ancillas are unentangled,
which basis-state testing cannot see. A stray diagonal gate would appear here as a phase, and phases
are exactly what Shor depends on.

\paragraph{The listings, too.}
Every listing in this part is extracted from \code{shor\_qiskit/ec\_*.py} by AST at build time, not
copied by hand, and the whole part is regenerated and diffed against the document by
\code{bench/ec\_check\_tex.py}. So a listing cannot drift from the function it claims to show:
change the code and the document either follows or the check fails. Where a listing is shortened,
the elision is visible as a bare \code{...} --- nothing is silently trimmed.

\paragraph{Property-based tests over random curves.}
Random prime, random non-singular $(a,b)$, random points and scalars. That is where curve-shape
assumptions hide: $a=0$, $b=0$, curves with a point of order two, primes where $2$ is or is not a
square. The harness shrinks failures --- staying inside each field's declared domain, and
re-checking its own result before reporting it --- and it has been validated against injected bugs,
where it isolates the exact offending modulus.

\begin{center}\small
\begin{tabular}{@{}ll@{}}
\toprule
what & result\\
\midrule
classical models (Kaliski, Montgomery, dialog, Jacobian) & exact\\
adders, comparators, modular arithmetic & exhaustive $p\le61$, random to $n=128$\\
Kaliski inversion and division & exhaustive; records and counter clean\\
affine point addition & every point pair, both control branches\\
projective point addition & every point pair; every $Z$ under \code{SHOR\_EC\_FULL}\\
in-place multiplication (dialog) & every $(x,y)$ for $p\le31$ ($p\le61$ full)\\
ECDLP end to end & every logarithm on two toy curves\\
\bottomrule
\end{tabular}
\end{center}

\subsection{The distribution, and the control that makes it mean something}\label{ec:sec:post}

On $y^2=x^3+5x+4\bmod7$ with $G=(0,5)$ of order $5$, and $Q=[2]G$, the measured peaks sit at
$(j_1,j_2)\in\{(0,0),(6,5),(2,3),(3,6),(5,2)\}$ carrying $20\%$, $10\%$, $10\%$, $10\%$, $10\%$ of
the shots. Rounding by $a_i=\mathrm{round}(j_i\cdot5/8)$ turns those into $(a_1,a_2) =
(1,2),(2,4),(3,1),(4,3)$ --- and every one of them satisfies $a_2=2a_1 \bmod 5$, which is
\eqref{ec:eq:post} with $k=2$.

%%LST ec_classical.ecdlp_postprocess doc=4%%

\begin{warnbox}[title={The control that catches a broken circuit}]
The correct $k$ must be the \emph{top} candidate, not merely \emph{a} candidate that verifies. On
both toy curves it is, for every logarithm, holding at least $30\%$ of the vote against a uniform
$20\%$ and $7.7\%$ respectively. And feeding the same post-processing a flat, random histogram
returns a top share of $21.2\%$ against a uniform $20\%$ --- no peak, no signal.

That control exists because an early version of this post-processing did not have it. Eker\aa{}'s
limited search \cite{ec:ekera} widens the rounding when $2^m$ is not a multiple of $r$, and with the
radius set too generously it enumerated the entire group: every $k$ ``verified'', on every input,
including noise. The bug was invisible in the pass/fail of an end-to-end test and obvious the
moment the vote shares were printed next to the uniform baseline.
\end{warnbox}

\section{What is not here}\label{ec:sec:gaps}

Named rather than glossed:

\begin{itemize}\setlength\itemsep{2pt}
\item \textbf{Register sharing} (\cite{ec:schrott26} Sec 3.1), which takes the dialog space from
$2.83n$ to $2.12n$ by shrinking $u$ and $v$ as the algorithm proceeds. It is left out not for being
probabilistic --- \S\ref{ec:sec:approx} builds and measures several circuits that are --- but
because its failure mode is a \emph{register overflow} rather than a wrong answer, so the
exact-simulation harness of \S\ref{ec:sec:verify} would be checking a circuit whose register layout
is itself only probably right. Priced in \S\ref{ec:sec:space}, not built.
\item \textbf{Algorithm 11} of \cite{ec:schrott26}, the pseudo-Mersenne controlled addition that
also handles $x+y=p$. That case cannot arise from random inputs, only in the first few
B\'ezout-replay iterations, and using the exact adder there is the same fix at negligible cost.
\item \textbf{Gidney's dirty-ancilla constant adder} \cite{ec:gidney25}. Constant addition here uses
clean ancillas: correct, and more qubits than the papers' accounting.
\item \textbf{The Ed25519 extension} of \cite{ec:kim26} Sec 4.4, in extended Edwards coordinates.
The curve model here is short Weierstrass only.
\item \textbf{Physical-level estimates} --- surface codes, QLDPC, wall-clock runtime. This part
stops at the logical layer, as Parts II--VI do.
\end{itemize}

\noindent
And one thing implemented in only one of its two forms: \cite{ec:kim26} distinguishes a
\emph{depth-optimised} Montgomery multiplier, which allocates fresh output qubits per two-operand
addition, from a \emph{qubit-optimised} one, which uncomputes and reuses them. The implementation
here always accumulates in place --- the qubit-optimised shape --- which is why the depth win from
the carry-save tree shows up smaller than the paper's depth-optimised figures.

\begin{keybox}[title={Where this part sits}]
The affine construction is \cite{ec:rnsl17} as improved by \cite{ec:hjnrs20}, with
\cite{ec:kim26}'s modifications: 2017--2020 with a 2026 finish. The projective addition and the
carry-save Montgomery multiplier are \cite{ec:kim26}. The Euclidean dialog, the approximate and
pseudo-Mersenne arithmetic, and the windowed addition are \cite{ec:schrott26}. What is deliberately
absent is the line that leaves this architecture entirely --- the residue-number-system
constructions that cut qubit counts by changing what a number \emph{is}, rather than how it is
added.

The honest limit of what runs on a laptop: the circuit with real arithmetic is verified on every
basis state, which for a permutation is the whole truth, but at $@@full_q@@$ qubits for a
seven-element field it cannot be run in superposition. The end-to-end demonstration uses a
permutation oracle of identical circuit shape. Both halves are tested; the join between them is an
argument, not a simulation.
\end{keybox}
